\documentclass[a4paper,12pt]{article}
\usepackage[utf8]{inputenc}
\usepackage[norsk]{babel}
\usepackage{hyperref}
\usepackage{geometry}
\geometry{margin=2.5cm}

\title{Krav til prosjektbeskrivelse (proposal) i LOG650\\Forskningsprosjekt: Logistikk og kunstig intelligens}
\date{}
\author{Høgskolen i Molde}

\begin{document}

\maketitle

Prosjektbeskrivelsen i LOG650 skal inneholde en foreløpig beskrivelse av prosjektet man vil gjennomføre. Siden den skrives i den innledende fasen av prosjektet, så er det naturlig at ikke alt er helt på plass, men noen momenter må være med. I dette dokumentet beskrives de ulike kravene for prosjektbeskrivelsen.

Merk at, siden prosjektbeskrivelsen er en tidlig versjon av prosjektet, så skal problemstilling, data, målfunksjon og avgrensninger beskrives kun med ord, og uten matematikk.

\section*{Innhold}

\begin{itemize}
    \item \textbf{Gruppemedlemmer:} Navnene på de som deltar i gruppen.
    \item \textbf{Område:} Hvilket fagområde prosjektet skal omhandle. Se oversikten over fagområder i kompendiet til emnet.
    \item \textbf{Bedrift (valgbart):} Hvilken bedrift man skriver i samarbeid med (hvis dette er avklart).
    \item \textbf{Problemstilling:} Hvilket problem man vil prøve å løse for bedriften. Hva man skal regne ut.
    \item \textbf{Data:} Hvilke data man vil trenge til prosjektet. Dersom reelle data er vanskelig å skaffe, kan det være aktuelt med simulerte data.
    \item \textbf{Målfunksjon:} Hvordan man måler forbedring i det man regner ut.
    \item \textbf{Avgrensninger:} Hva som ikke skal være med i modellen eller analysen.
\end{itemize}

\section*{Øvrige krav}

\begin{itemize}
    \item Prosjektbeskrivelsen skal være på maksimalt 2 sider.
    \item Problemstillingen må være kvantitativ, altså regnemessig.
\end{itemize}

\end{document}
